\documentclass{standalone}
\usepackage{circuitikz}
\usetikzlibrary{calc}
\definecolor{circuitnotemark}{HTML}{FE00FE}
\begin{document}
\begin{circuitikz}[american, line width=0.8pt]
\ctikzset{bipoles/length=1.2cm}
\foreach \x in {0,2,4,6,8,10,12,14} {\foreach \y in {0,-2,-4} {\fill[gray, opacity=0.35] (\x,\y) circle (1.1pt);}}
\node[gray, font=\scriptsize] at (0,0.84) {1};
\node[gray, font=\scriptsize] at (2,0.84) {2};
\node[gray, font=\scriptsize] at (4,0.84) {3};
\node[gray, font=\scriptsize] at (6,0.84) {4};
\node[gray, font=\scriptsize] at (8,0.84) {5};
\node[gray, font=\scriptsize] at (10,0.84) {6};
\node[gray, font=\scriptsize] at (12,0.84) {7};
\node[gray, font=\scriptsize] at (14,0.84) {8};
\node[gray, font=\scriptsize, anchor=east] at (-0.84,0) {a};
\node[gray, font=\scriptsize, anchor=east] at (-0.84,-2) {b};
\node[gray, font=\scriptsize, anchor=east] at (-0.84,-4) {c};
\coordinate (a1) at (0,0);
\coordinate (a2) at (2,0);
\coordinate (a4) at (6,0);
\coordinate (a5) at (8,0);
\coordinate (a7) at (12,0);
\coordinate (a8) at (14,0);
\coordinate (c1) at (0,-4);
\coordinate (c2) at (2,-4);
\draw (a1) to[R, l_=$R_{1}$] (a2); % line 3
\draw (a4) to[R, l_=$R_{load}$] (a5); % line 4
\draw (a7) to[R, l_=$R$] (a8); % line 5
\draw (c1) to[esource, l_=$V_{cc2}$] (c2); % line 6
\draw (0.73,-4) -- ++(0.2,0); % line 6
\draw (0.83,-4.1) -- ++(0,0.2); % line 6
\draw (1.07,-4) -- ++(0.2,0); % line 6
\path (-1.286,-1.993) rectangle (3.286,-0.807); % line 8
\node[anchor=west, inner sep=0, circuitnotemark, font=\large] at (1,-1.4) {X}; % line 8
\path (4.333,-1.993) rectangle (9.667,-0.807); % line 9
\node[anchor=west, inner sep=0, circuitnotemark, font=\large] at (7,-1.4) {X}; % line 9
\path (10.841,-1.993) rectangle (15.159,-0.807); % line 10
\node[anchor=west, inner sep=0, circuitnotemark, font=\large] at (13,-1.4) {X}; % line 10
\path (-1.413,-5.993) rectangle (3.413,-4.807); % line 11
\node[anchor=west, inner sep=0, circuitnotemark, font=\large] at (1,-5.4) {X}; % line 11
\path (18,-7.895) rectangle (31.911,0.593); % line 12
\node[anchor=west, inner sep=0, circuitnotemark, font=\large] at (18,0) {X}; % line 12
\node[anchor=west, inner sep=0, circuitnotemark, font=\large] at (18,-0.487) {X}; % line 12
\node[anchor=west, inner sep=0, circuitnotemark, font=\large] at (18,-0.974) {X}; % line 12
\node[anchor=west, inner sep=0, circuitnotemark, font=\large] at (18,-1.46) {X}; % line 12
\node[anchor=west, inner sep=0, circuitnotemark, font=\large] at (18,-1.947) {X}; % line 12
\node[anchor=west, inner sep=0, circuitnotemark, font=\large] at (18,-2.434) {X}; % line 12
\node[anchor=west, inner sep=0, circuitnotemark, font=\large] at (18,-2.921) {X}; % line 12
\node[anchor=west, inner sep=0, circuitnotemark, font=\large] at (18,-3.408) {X}; % line 12
\node[anchor=west, inner sep=0, circuitnotemark, font=\large] at (18,-3.895) {X}; % line 12
\node[anchor=west, inner sep=0, circuitnotemark, font=\large] at (18,-4.381) {X}; % line 12
\node[anchor=west, inner sep=0, circuitnotemark, font=\large] at (18,-4.868) {X}; % line 12
\node[anchor=west, inner sep=0, circuitnotemark, font=\large] at (18,-5.355) {X}; % line 12
\node[anchor=west, inner sep=0, circuitnotemark, font=\large] at (18,-5.842) {X}; % line 12
\node[anchor=west, inner sep=0, circuitnotemark, font=\large] at (18,-6.329) {X}; % line 12
\node[anchor=west, inner sep=0, circuitnotemark, font=\large] at (18,-6.816) {X}; % line 12
\node[anchor=west, inner sep=0, circuitnotemark, font=\large] at (18,-7.302) {X}; % line 12
\node[anchor=south west, inner sep=2pt, circuitnotemark, font=\LARGE\bfseries] (circuittitle) at (current bounding box.north west) {X};
\path (circuittitle.south west) rectangle ++(6.035,0.853);
\end{circuitikz}
\end{document}
